% Journal options: aps + prl/pra/prb/prc/prd/pre/prx for the Physical Review family,
% or aip + jcp/apl/jap/... for AIP Publishing. "reprint" gives the published two-column
% look, "preprint" a one-column manuscript.
\documentclass[aps,pra,reprint,amsmath,amssymb,superscriptaddress]{revtex4-2}
\usepackage{graphicx}
\usepackage{bm}
\usepackage[colorlinks=true,allcolors=blue]{hyperref}

\begin{document}

\title{Title of the Letter or Article}

\author{First Author}
\email[Correspondence: ]{first.author@example.edu}
\affiliation{Department of Physics, Example University, City, Country}

\author{Second Author}
\affiliation{Example Institute, City, Country}

\date{\today}

\begin{abstract}
A single paragraph that states the problem, the method and the main quantitative
result, understandable without the rest of the paper. No citations.
\end{abstract}

\maketitle

\section{Introduction}
Introduce the physical question and cite earlier work~\cite{einstein35}.
Equations use the usual environments and are numbered automatically:
\begin{equation}
  i\hbar \frac{\partial}{\partial t}\Psi(\mathbf{r},t) = \hat{H}\,\Psi(\mathbf{r},t).
  \label{eq:schroedinger}
\end{equation}
Refer to them as Eq.~(\ref{eq:schroedinger}).

\section{Model}
Define the system and the approximations. Wide equations can be broken up with
\texttt{align}, or set across both columns with the \texttt{widetext} environment.

\section{Results}
% Uncomment and add the file: Fig.~\ref{fig:setup} then works.
% \begin{figure}
%   \includegraphics[width=\columnwidth]{setup.pdf}
%   \caption{Caption below the figure.}
%   \label{fig:setup}
% \end{figure}
Describe your results here, with tables in the \texttt{table} environment
(\texttt{ruledtabular} gives the journal rules).

\begin{table}
\caption{A small example table.}
\begin{ruledtabular}
\begin{tabular}{lcc}
 Quantity & Theory & Experiment \\
 \hline
 $\omega_0/2\pi$ (MHz) & 5.02 & 4.98(3) \\
 $Q$ & $1.2\times10^{4}$ & $1.1(1)\times10^{4}$ \\
\end{tabular}
\end{ruledtabular}
\end{table}

\section{Conclusion}
Summarise what was shown and what follows.

\begin{acknowledgments}
We thank colleagues and funding agencies here.
\end{acknowledgments}

\bibliography{refs}

\end{document}
